\subsection{An exact exclusion of an early parameter family}
\begin{finding}{P-21}{the exceptional projective-line branch at a support-thirteen ten-plane}
Work over an algebraically closed field of characteristic zero.
Let $T$ be the determinant residual tensor for
\[
 S=(0,1,2,3,6,7,9,12,18,19),
 \qquad \dim A=10,\quad\dim B=\dim C=13.
\]
Use the quotient $V=(A\otimes B)/T(C^*)$ and its unit-pivot
weight bases specified in P-17. Let $F_0$ be the sum of the
one-dimensional weight blocks $56$ and $76$, and let $K$
be the two-dimensional block $48$. There is no $H$-invariant
seventeen-plane $E\supset T(C^*)$ satisfying
\[
 \rank\Psi_B(E)\leq204
 \tag{P21.1}
\]
whose quotient $F=E/T(C^*)$ has an invariant flag containing
$F_0\subset F_0+L$, for some line $L\subset K$.
Thus this entire early projective-line branch in the
ten-plane test is excluded. Other extension spaces for
this subset are not excluded by this statement alone.
\end{finding}
\begin{proof}
All quotient block numbers in this proof are reconstructed
from $S$ by the following unambiguous rule. Order pairs
$(q,v)$ lexicographically, where $q\in S$ and $v$ ranges
over the used variables in increasing order. Group them
by weight $\operatorname{wt}(v)-\operatorname{wt}(q)$,
ordering groups by their first occurrence. In every
image weight of $U=T(C^*)$, eliminate the last coordinate
using its unit coefficient in $U$. Discard any resulting
zero-dimensional group and number the remaining groups
from zero. The coordinates of $K$ are the classes of
$(q,v)=(7,0),(9,1)$, while $F_0$ has coordinates
$(9,0)$ and $(19,0)$. Here $v=4r+c$ denotes a matrix
entry, and $q$ is the rectangle index used throughout.

The six adjacent root matrices on $V$ are those of
P-17. The standalone verifier below reconstructs them
and their integral equivariance identities. Direct exact
kernel calculations show that the common root kernel
on $V/F_0$ is the direct sum of $K$ and one line at each
of the eight weights
\[
 23,27,32,58,61,63,70,77.
 \tag{P21.2}
\]
This also verifies that $F_0$ itself is invariant.

Suppose first that the last flag line is a line already
annihilated by the roots on $V/F_0$, at a weight in
(P21.2). Denote that fixed line by $X_i$. The possible
final space then has the form
$F_0+X_i+L$ with $L\subset K$ arbitrary. Write
$L=\langle\lambda e_1+\mu e_2\rangle$ in the displayed
quotient basis. The $B$-exterior map of the fixed
sixteen-dimensional space $U+F_0+X_i$ has rank at least
$195$, certified by a nonzero order-$195$ minor.
The new line contributes thirteen possible image
columns, in distinct weights
$\operatorname{wt}(K)+\operatorname{wt}(b_j)$.

For each such weight the checker computes integer
annihilators of the old image block. Applying those
annihilators to the two new coefficient columns gives
an explicitly recorded two-column matrix $C_j$.
If $C_j(\lambda,\mu)^{\mathsf t}\ne0$, that new column
is outside the old image. The distinct weights make
these increments independent. A rank-two $C_j$ never
vanishes on a projective point; a rank-one $C_j$ vanishes
at its one kernel point; and a zero $C_j$ contributes
no increment. This exactly reduces the lower-bound
calculation to the generic point and a finite list
of explicitly calculated kernel points. The resulting
minimum numbers of nonzero new columns are
\[
\begin{array}{c|rrrrrrrr}
i&23&27&32&58&61&63&70&77\\\hline
\min\#\{j:C_j(\lambda,\mu)^{\mathsf t}\ne0\}
 &11&11&11&10&11&11&11&10.
\end{array}
\]
Every member of these eight families therefore has
$B$-exterior rank at least $195+10=205$, contradicting
(P21.1). If the last flag line has the same weight
as $K$, the final space contains all of $K$; a separate
nonzero order-$205$ minor excludes this fixed plane.

It remains to treat last lines which are not in the
old common kernel. A root applied to such a line must
have nonzero image in $L$. Its weight therefore has
a root arrow into the weight of $K$. There is at most
one such adjacent root for a given source weight.
Impose all the other root-kernel equations first.
The complete list of nonzero possibilities is
\[
\begin{array}{c|c|c|c}
\text{source block}&\text{preimage vector}
 &\text{root image in }K&\text{required }[\lambda:\mu]\\\hline
13&(1,1)&(-1,-1)&[1:1]\\
38&(0,1)&(0,1)&[0:1]\\
49&(1,1)&(0,-2)&[0:1]\\
51&(1,1)&(2,2)&[1:1].
\end{array}
\]
Each preimage solution space is one-dimensional and
has no vector with zero root image except zero. The
checker forms every root arrow, so this list is
exhaustive. Consequently a genuinely new last line
is possible only at $L=\langle(1,1)\rangle$ or
$L=\langle(0,1)\rangle$.

For each of these two fixed lines, recompute the
common root kernel in $V/(F_0+L)$. It consists of
eleven distinct one-dimensional weight spaces; no
weight has a higher-dimensional kernel. Thus there
are exactly eleven possible next lines for each
special point. Each of the resulting twenty-two
seventeen-planes has a certified nonzero order-$205$
$B$-exterior minor. This excludes the remaining cases.

The self-contained program
\path{verify_early_borel_branch.py} uses only NumPy
and the Python standard library. It reconstructs
$T$, the quotient gauge, all six root matrices,
the exact rational kernels, all old-image annihilators,
the matrices $C_j$, and every exceptional kernel point.
It independently evaluates the eight baseline minors,
the full-$K$ minor and the twenty-two special-point
minors modulo the explicitly verified prime $1000003$.
The thirty-one nonzero residues certify nonzero integer
minors in characteristic zero. The complete data,
including all annihilator identities and minor indices,
are written to \path{early_borel_branch_certificate.json}.
No sampled parameter value is substituted for a
projective family in this argument.
\end{proof}

\paragraph{Dictionary and scope.}
P-21 refines the proper Grassmannian incidence of P-17
by a flag incidence and its projective-line fiber.
The root maps and weight summands are maps of the
associated equivariant vector bundles. The argument
retains the actual maps, their kernels and their
integer minors; their classes in $K_0$ alone would
not supply the exclusion. The applicable EGA~II,
SGA~3 and FGA locators are those in the common dictionary.
The conclusion closes one concrete parameter branch,
and is not a uniform assertion that $\mu_{10}\geq18$.
